\section{Conclusiones}
El objetivo retrospectivo se respondió bajo condiciones declaradas: las masas de reposo y demás propiedades del inventario HRG, combinadas con el exceso bariónico observado y una temperatura elegida, permitieron reconstruir poblaciones separadas de bariones y antibariones. A 150 MeV se infirieron aproximadamente $1.7222\times10^{37}$ cm$^{-3}$ de cada lado, con exceso de $4.0110\times10^{29}$ cm$^{-3}$. A 40 MeV se distinguieron $1.9166\times10^{28}$ bariones y $1.4894\times10^{28}$ antibariones por cm$^3$. El rendimiento total bajó 7.0323 órdenes de magnitud.

La abundancia actual sola no recuperó una población inicial única: el mapa del exceso tuvo rango uno y conservó cualquier adición igual de partículas y antipartículas. La condición de equilibrio a temperatura fijada fue información adicional necesaria. El contraste con el límite nucleónico publicado cuantificó una participación de otros bariones de 74.42\% a 150 MeV y 1.525\% a 40 MeV, coherente con su supresión térmica.

La referencia publicada $T_c=156.5\pm1.5$ MeV fundamentó una comprobación de la temperatura de partida. A su valor central, el mismo HRG dio $2.6575\times10^{37}$ cm$^{-3}$ por lado al redondeo mostrado, con un total 54.31\% mayor que a 150 MeV. Esta sensibilidad cuantifica la importancia del punto inicial, sin identificarlo con un inicio de aniquilación medido.

Los resultados son una reconstrucción reproducible condicionada, sin demostrar un descubrimiento físico nuevo o un historial único de aniquilación. Las densidades actuales de referencia, $n_{b0}=(2.514\pm0.017)\times10^{-7}$ cm$^{-3}$ y el equivalente primordial de H de $(1.898\pm0.015)\times10^{-7}$ cm$^{-3}$, proceden de conversiones de entradas observacionales. Determinar eventos acumulados, inicio cinético y abundancias nucleares finales requiere tasas de reacción y evolución cosmológica.

\section*{Disponibilidad de datos y código}
El suplemento incluye catálogos fijados y SHA-256; densidades numéricas y de entropía; 3544 registros por especie de partícula y temperatura para el control PDG2025; 24 estados de inferencia inversa y 12 evaluaciones de temperatura inicial; el contraste nucleónico; la prueba de rango; las conversiones cosmológicas; controles numéricos; cinco figuras por idioma y fuentes del manuscrito. Los cálculos se ejecutaron con la implementación Python independiente incluida, utilizando los datos de Thermal-FIST y conservando su licencia original.

\section*{Asistencia en la preparación}
Se utilizó asistencia de inteligencia artificial para organizar la redacción y preparar código, figuras y maquetación. Las sumas, integrales y soluciones inversas descritas se ejecutaron y sus resultados se conservaron. La evaluación científica y la decisión de publicación corresponden al autor.
